\documentclass{article}
\usepackage{amsmath,amssymb}
\newcommand{\Pn}{P_n}
\newcommand{\e}{\mathrm{e}}
\DeclareMathOperator*{\argmax}{arg\,max}
\title{The Cutoff Rule in Sequential Selection: Exact Finite-Sample Optimum and Simulation}
\begin{document}
\begin{abstract}
A decision maker sees $n$ candidates one by one in random order, must accept or reject each on the spot,
and wants the best one. The cutoff rule skips the first $r-1$ candidates and then takes the first who beats
all predecessors. We restate the exact success probability, derive the large-$n$ optimum $r/n \to 1/\e$, and
compare both with a simulation.
\end{abstract}
\section{Model}\label{sec:model}
Candidates arrive in uniformly random order and can be ranked without ties. The rule with cutoff $r$
succeeds with probability
\begin{equation}\label{eq:prob}
  \Pn(r) = \frac{r-1}{n} \sum_{k=r}^{n} \frac{1}{k-1}, \qquad r \ge 2 .
\end{equation}
\section{Optimal cutoff}\label{sec:opt}
With $x = r/n$ held fixed as $n \to \infty$,
\begin{equation}\label{eq:limit}
  P(x) = -x \ln x, \qquad x^\star = \argmax_x P(x) = 1/\e, \qquad P(x^\star) = 1/\e \approx 0.368 .
\end{equation}
\section{Simulation}\label{sec:sim}
We simulate $20{,}000$ arrival orders per cutoff for $n = 20$ and $n = 100$.
\end{document}
